\section{Robustness and Measurement Scope}
\label{sec:robustness}

\subsection{Candidate Allocation}

Table~\ref{tab:geographic-agreement} reports all three allocation rules already
computed for the geographic comparison. The island--coastal difference is
0.461 under equal shares, 0.631 under projection-score weighting, and 0.715
under Top-1 allocation. The ordering is stable across these rules, while the
coefficient magnitudes are not. The paper therefore reports the allocation
rule with each result.

\subsection{Landing-Region Resolution}

The existing landing-region-diameter analysis covers the A-Root subset, not all
18 measurements. It varies maximum diameter over 10, 20, 30, 40, and 50~km
while holding the 50-km endpoint catchment and 5-ms RTT tolerance fixed.
Table~\ref{tab:diameter-sensitivity} reports the recorded values. Exact
landing-pair and cable candidate sets have mean Jaccard 1.0 on segments shared
with the 30-km baseline, but the candidate-bearing segment population and the
number of auditable units change with diameter.

\begin{table}[t]
  \centering
  \caption{A-Root landing-region diameter sensitivity.}
  \label{tab:diameter-sensitivity}
  \footnotesize
  \setlength{\tabcolsep}{3.2pt}
  \begin{tabular}{rrrr}
    \toprule
    Diameter & Single-corridor & Median Top-2 & Units \\
    \midrule
    10 km & 47.4\% & 53.6\% & 16 \\
    20 km & 46.2\% & 62.8\% & 15 \\
    30 km & 54.8\% & 64.4\% & 15 \\
    40 km & 90.1\% & 81.6\% & 13 \\
    50 km & 94.8\% & 84.9\% & 13 \\
    \bottomrule
  \end{tabular}
\end{table}

The 10--30-km settings retain a finer coastal partition than the 40--50-km
settings. Because the existing analysis is limited to A-Root and does not
recompute the main geographic Layer Agreement result, it does not establish
that the headline geographic difference is invariant to landing-region
diameter.

\subsection{Dependence on Geolocation and Candidate Construction}

Every corridor projection begins with hop geolocation. A hop placed farther
inland can remove a nearby landing candidate; a location shifted toward a coast
can add one. The 50-km catchment is the fixed setting used in the reported
analysis. No location-perturbation or alternative-geolocation result is
available in the existing experiment set, so the effect of residual IPinfo
error is not bounded here.

The projection also lacks physical-route ground truth. Geographic, lifecycle,
and propagation constraints define feasible candidates rather than observed
cable use. Exact landing-pair, cable, and corridor identities remain linked in
the pipeline outputs, but no completed manual corridor audit is reported in
this draft.

\subsection{Snapshot and Population Scope}

All 18 measurements cover 00:00--01:00 UTC on July~1, 2026. Routing, service
deployment, and cable availability can change outside this window, so the
reported distributions describe one aligned snapshot. RIPE Atlas probe density
varies across countries and networks, and all quantities are observation-
weighted over participating probes.

The service measurements and dynamic multi-target references have different
target populations. Family comparisons describe those observed populations
and are not matched estimates of a service-selection or deployment effect.
Anycast observations jointly reflect deployed footprint, BGP selection,
interconnection, and source-network context.

Cable lifecycle filtering uses the inventory fields available for the
measurement date. These fields do not report transient failures or outages,
and the cable inventory can be revised. The results therefore measure feasible
candidate support under the recorded metadata, not operational redundancy,
hazard correlation, rerouting capacity, or repair performance.
